\section{Métodos alternativos de estimação do SSF com tempo de trânsito}
\label{sec:modelos_alternativos_est-ssf}

Até o momento, foram tratados métodos para estimação do SSF que ou são a \ingles{Travel-Time Tomography} tradicional (\cref{sec:modelo-padrão-TTT}), ou são modificações pequenas em cima da TTT tradicional, incluindo informações diferentes (\cref{sec:modelos_adaptados_ttt}), ou com um método diferente para o passo da inversão (\cref{sec:svd}).

Dado isso, é possível utilizar a estrutura da TTT (de minimização do erro de tempo de trânsito entre o modelo e o resultado), e aliá-la a outras ferramentas de modelagem, de forma a obter um resultado que mais proximamente se alinhe ao SSF real. Isso foi tentado através da SVD anteriormente (\cref{sec:svd}), mas como mostrado na \cref{sec:simulation_and_results} isso não foi suficiente para alcançar resultados satisfatórios na situação em camadas; essa sendo a mais alinhada com imposições práticas.

\subsection{Método por parametrização de Munk}
\label{subsec:sec4:metodo_parametrizado_por_munk}

O primeiro método é uma adaptação do proposto por \citet{wu_inversion_2023}, expandindo-o para o caso multi-dimensional, e com uma escolha de parâmetros levemente modificada.

Para esse método, retornaremos ao perfil de Munk da \cref{eq:sec4:perfil_de_munk}. Consideraremos um ambiente bidimensional (sem perda de generalidade), com $\sz{N_{\z}}{N_{\x}}$ células. Cada coluna de células (com posição $x$) será regida por um perfil de Munk diferente, com parâmetros $v_x$, $\epsilon_x$, $z_x$, e $B_x$; e, disso, pela \cref{eq:sec4:perfil_de_munk},
\begin{equation}
	v(x,z) = V_x \pts{ 1 + \epsilon_x \bts{ e^{2\frac{z-z_x}{B_x} } + 2\frac{z-z_x}{B_x} - 1 } }
	\label{eq:sec5:perfil_munk_parametrizado}
\end{equation}

Denota-se $\bvV,\bvep,\bvz,\bvB \in \re^{\sz{N_{\x}}{1}}$ como vetores dos parâmetros de Munk, e $\bvbeta \in \re^{\sz{4N_{\x}}{1}}$ como um vetor de parâmetros, contemplando todos os $4N_{\x}$ parâmetros de modelagem, dado por
\begin{equation}
	\bvbeta = \tr{\bts{\bvV[T] ~ \bvep[T] ~ \bvz[T] ~ \bvB[T]}}
	\label{eq:sec5:definicao_vetor_bvbeta}
\end{equation}

Por fim, $\bvv(\bvbeta)$ é o SSF para esse conjunto de parâmetros. Ou ainda, o campo de vagarosidade (também denotado SSF, por abuso de notação)
\begin{equation}
	\bvs(\bvbeta) = \frac{1}{\bvv(\bvbeta)}
	\label{eq:sec5:ssf_usando_bvbeta}
\end{equation}

\subsubsection*{Função-custo e erro quadrático}

De forma semelhante à TTT, será usada uma minimização iterativa, alternando entre encontrar o SSF (agora parametrizado por $\bvbeta$) e a matriz de distâncias. Matematicamente, $\bvbeta{i}$ é o vetor de parâmetros da $i$-ésima iteração, e $\bvR{i} \equiv \bvR(\bvs(\bvbeta{i})) \in \re^{\sz{M}{N}}$ é a matriz de distâncias para o SSF definido por $\bvbeta$. A função custo será
\begin{equation}
	E(\bvbeta{i}) = \norm{\bvR{i-1} \bvs(\bvbeta{i}) - \bvt{\obs}}^2
	\label{eq:sec5:funcao-custo-beta}
\end{equation}
onde assume-se $\bvR{i-1}$ constante. Expandindo a norma $\ell_2$ da \cref{eq:sec5:funcao-custo-beta}, tem-se
\begin{equation}
	E(\bvbeta{i}) = \sum_{m=0}^{M-1} E_m^2(\bvbeta{i})
\end{equation}
onde
\begin{equation}
	E_m(\bvbeta{i}) = \sum_{n=0}^{N-1} r_{m,n;i-1} s_{n;i} - t_{\obs;m}
\end{equation}
em que $r_{m,n;i-1} = \el{\bvR{i-1}}{m,n}$ é a distância percorrida pela $m$-ésima onda (fonte-receptor) dentro da $n$-ésima célula na $i-1$-ésima iteração; $s_{n;i}$ é a vagarosidade da $n$-ésima célula na $(i)$-ésima iteração (a atual), sendo função de $\bvbeta{i}$ (ou mais especificamente, dos elementos de $\bvbeta{i}$ que correspondem à coluna na qual a $n$-ésima célula se encontra); e $t_{\obs;m}$ é o tempo de trânsito da $m$-ésima onda.

Denota-se também $\bve{i-1} \in \re^{\sz{M}{1}}$ um vetor dos resíduos, com $\el{\bve{i}}{m} = E_m(\bvbeta{i})$.

Ressaltaremos que, nas formulações anteriores, foi usada uma notação em $i$ que definia a próxima iteração ($i+1$) com base na atual. Por limpeza e espaço, aqui será usada uma notação que define a iteração atual $i$ com base na anterior $i-1$.

\subsubsection*{Minimização e o algoritmo de Gauss-Newton}

É notável que a relação entre a função-custo $E(\bvbeta{i})$ e a variável de minimização $\bvbeta{i}$ é não-linear. Enquanto na TTT clássica, a função-custo é estritamente uma função quadrática em relação aos parâmetros de otimização (o vetor de vagarosidade $\bvs{i}$), aqui o comportamento de $\bvs{i}$ é regido pelos parâmetros de $\bvbeta{i}$. Isso faz com que a solução por inversão não seja direta, e tampouco possua uma forma fechada para sua solução.

Para tratar desse problema não-linear, utilizaremos o algoritmo de Gauss-Newton (GNa). Para tal, denota-se $\bvbeta{i,j}$ como a $j$-ésima iteração (dentro do GNa) do vetor de parâmetros, na $i$-ésima iteração da estimação do SSF. Também escrevemos $\bvJ{i,j} \in \re^{\sz{M}{4N_{\x}}}$ como o Jacobiano de $\bve{i,j}$ em relação a $\bvbeta{i,j}$; ou ainda,
\begin{equation}
	\el{\bvJ{i,j}}{m,k} = \frac{\partial E_m}{\partial \beta_{k;i,j}}(\bvbeta{i,j})
\end{equation}

O GNa é inicializado com $\bvbeta{i,\{j=0\}} = \bvbeta[o]{i-1}$, recebendo o vetor de parâmetros ótimo da solução anterior (em $i$).

Lembrando que $s_{n;i-1,j} = \frac{1}{v_{n;i-1,j}}$, é fácil mostrar que
\begin{subequations}
	\begin{align}
		\pdv{E_m}{V_{k;i,j}} & = \sum_{n=0}^{N-1} f_{m,n;i,j} \pts{1 + \epsilon_{k;i,j} g_z(z_{k;i,j},B_{k;i,j}) } \\
		\pdv{E_m}{\epsilon_{k;i,j}} & = V_{k;i,j} \sum_{n=0}^{N-1} f_{m,n;i,j} g_z(z_{k;i,j},B_{k;i,j}) \\
		\pdv{E_m}{z_{k;i,j}} & = V_{k;i,j} \epsilon_{k;i,j} \frac{2}{B_{k;i,j}} \sum_{n=0}^{N-1} f_{m,n;i,j} \pts{e^{-\eta_z(z_{k;i,j},B_{k;i,j})} - 1} \\
		\pdv{E_m}{B_{k;i,j}} & = V_{k;i,j} \epsilon_{k;i,j} \sum_{n=0}^{N-1} f_{m,n;i,j} \frac{\eta_z(z_{k;i,j},B_{k;i,j})}{B_{k;i,j}} \pts{e^{-\eta_z(z_{k;i,j},B_{k;i,j})} - 1}
	\end{align}
	\begin{gather}
		f_{m,n;i,j} = - \Delta_{n,k} s^2_{n;i,j} r_{m,n;i-1-1} \\
		g_z(z_{k;i,j},B_{k;i,j}) = \bts{e^{-\eta_z(z_{k;i,j},B_{k;i,j})} + \eta_z(z_{k;i,j},z_{k;i,j}) - 1 } \\
	\eta_z(z_{k;i,j},B_{k;i,j}) = \frac{2\pts{z_n - z_{k;i,j}}}{B_{k;i,j}}
	\end{gather}	
\end{subequations}
em que $V_{k;i,j}$, $\epsilon_{k;i,j}$, $z_{k;i,j}$ e $B_{k;i,j}$ são os $[k,N_{\x}+k,2N_{\x}+k,3N_{\x}+k]$-ésimos elementos de $\bvbeta{i,j}$, e $z_n$ é a profundidade da $n$-ésima célula. $\Delta_{n,k}$ é uma versão modificada do delta de Dirac, sendo $1$ se a $n$-ésima célula pertence à $k$-ésima coluna, e $0$ caso contrário.

Tendo o Jacobiano $\bvJ{i,j}$, o vetor de resíduos $\bve{i,j}$, e o vetor de parâmetros $\bvbeta{i,j}$, pelo GNa tem-se então
\begin{equation}
	\bvbeta{i,j+1} = \bvbeta{i,j} - \inv*{\bvJ[T]{i,j} \bvJ{i,j}} \bvJ[T]{i,j} \bve{i,j}
\end{equation}

Esse processo iterativo é repetido em $j$ até ocorrer a convergência em $\bvbeta{i,j+1}$; momento no qual se extrai $\bvbeta[o]{i} = \bvbeta{i,j}$.

Para que esse processo seja válido, é necessário e suficiente que $\rank{\bvJ{i,j}} = 4N_{\x} \leq M$. Contudo, essa condição não é garantida verdadeira.

\subsubsection*{Regularização com referência}

O perfil de Munk é definido com um conjunto de parâmetros, geralmente $V_0 = 1500\si{\meter/\second}$, $\epsilon_0 = 0.00737$, $z_0 = 1300\si{\meter}$, e $B_0=1300\si{\meter}$. Esses valores podem ser usados como referências para a minimização apresentada anteriormente, através de uma regularização aditiva.

Definimos $F(\bvbeta{i,j})$ como um erro quadrático entre os parâmetros a sere estimados e os valores-base de referência, dado por
\begin{equation}
	F(\bvbeta{i,j}) = \norm{\bvA \pts{\bvbeta{i,j} - \bvbeta{\rfr}}}^2
\end{equation}
em que $\bvbeta{\rfr} \in \re^{\sz{4N_{\x}}{1}}$ é o vetor de referência, em que cada bloco de $N_{\x}$ valores é igual ao respectivo valor de referência citado anteriormente; e $\bvA \in \re^{\sz{4N_{\x}}{4N_{\x}}}$ é uma matriz diagonal com parâmetros de regularização, dada por
\begin{equation}
	\bvA = \begin{bmatrix}
		\alpha_v \Id & \bv0 & \bv0 & \bv0 \\
		\bv0 & \alpha_{\epsilon} \Id & \bv0 & \bv0 \\
		\bv0 & \bv0 & \alpha_z \Id & \bv0 \\
		\bv0 & \bv0 & \bv0 & \alpha_B \Id
	\end{bmatrix}
	\label{eq:s1:matriz_param-regularizacao}
\end{equation}

Os elementos de $\bvA$ servem simultaneamente para controlar o quão ``pesada'' é a regularização, e também normalizar o impacto dos parâmetros, dada a diferença de ordem entre eles.

Por usarmos uma regularização quadrática, seu Jacobiano será $\bvF{i,j}' = \bvA$. Com isso, a iteração do GNa se torna
\begin{equation}
	\bvbeta{i,j+1} = \bvbeta{i,j} - \inv*{ \bvJ[T]{i,j} \bvJ{i,j} + \bvF[pT]{i,j} \bvF[p]{i,j} } \pts{ \bvJ[T]{i,j} \bve{i,j} + \bvF[pT]{i,j} \bvF[p]{i,j} \bts{ \bvbeta{i,j} - \bvbeta{\rfr} } }
\end{equation}

Alternativamente, essa regularização poderia ser feita sobre o SSF reconstruído usando as \cref{eq:sec5:perfil_munk_parametrizado,eq:sec5:definicao_vetor_bvbeta,eq:sec5:ssf_usando_bvbeta}, ao invés do vetor de parâmetros; ou uma regularização que relacionasse os parâmetros entre colunas adjacentes. A regularização utilizada foi escolhida por simplicidade matemática, embora as duas outras avenidas permaneçam em aberto como vias de exploração futura.

\subsection{Método por EOFs}
\label{subsec:sec4:metodo_por_eof}

Esse método será semelhante ao proposto por \cite{radhakrishnan_inversion_2024}, com as devidas modificações e adaptações para o caso multi-dimensional, além de algumas outras alterações.

A modelagem por funções empíricas ortogonais (EOFs, do inglês \ingles{Empirical Orthogonal Functions}) são uma abordagem possível para reconstrução do SSF. Nela, cada coluna do SSF não é estimada diretamente, mas como uma combinação linear dentre possíveis perfis previamente estimados, a partir de um banco de dados.

\subsubsection*{Construção da base de EOFs}

A partir de um banco de dados com $K$ medições de temperatura, pressão, e salinidade (TPS) ao longo da coluna d'água em uma localidade ou região, é possível determinar o SSP (unidimensional) nessa localidade através de formulações já estabelecidas \cite{delgrosso_new_1974,mackenzie_nineterm_1981,wilson_equation_1960}.

Denota-se $\bvs{k} \in \re^{\sz{N_{\z}}{1}}$ o SSP para o $k$-ésimo registro, e $\bvs[b]$ o perfil médio entre os $K$ registros,
\begin{equation}
	\bvs[b] = \frac{1}{K} \sum_{k=0}^{K-1} \bvs{k}
\end{equation}
a partir do qual define-se a matriz de covariância entre os $K$ registros $\Corr{\bvs} \in \re^{\sz{N_{\z}}{N_{\z}}}$,
\begin{equation}
	\Corr{\bvs} = \frac{1}{K} \sum_{k=0}^{K-1} \pts{\bvs{k} - \bvs[b]} \tr{\pts{\bvs{k} - \bvs[b]}}
\end{equation}

Assumindo que $\Corr{\bvs}$ seja uma matriz de posto cheio\footnote{Isso pressupõe que $K > N_{\z}$, indicando mais medições que colunas, o que pode ser inviável ou indesejado na prática.}, com $\rank{\Corr{\bvs}} = N_{\z}$, então $\kernel{\Corr{\bvs}} = \re^{\sz{N_{\z}}{1}}$. Isso significa que qualquer vetor $\bvs[t]$ (inclusive SSP fisicamente impossíveis) pode ser obtido como uma combinação linear das colunas de $\Corr{\bvs}$. Por simplicidade, escreve-se que
\begin{equation}
	\bvs[t] = \bvs[b] + \Corr{\bvs} \bvalpha
	\label{eq:sec5:ssp_combinacao-linear}
\end{equation}
onde $\bvalpha$ é um vetor de coeficientes. Essa explicitação da vagarosidade média $\bvs[b]$ é feita por $\Corr{\bvs}$ considerar somente a variação dos vetores $\bvs{k}$ em relação à média, e (em geral) vai fazer com que os coeficientes de $\bvalpha$ sejam mais bem comportados.

É possível decompor $\Corr{\bvs}$ em função de seus autovalores e autovetores, levando a
\begin{equation}
	\begin{split}
		\Corr{\bvs}
		& = \sum_{n=0}^{N_{\z}-1} \lambda_n \bvv{n} \bvv[T]{n} \\
		& = \bvV \bvLa \inv{\bvV}
	\end{split}
\end{equation}
com $\abs{\lambda_1} \geq \abs{\lambda_2} \geq \cdots \geq \abs{\lambda_{N_{\z}-1}}$ os autovalores de $\Corr{\bvs}$, em ordem decrescente; e $\bvv{n}$ a $n$-ésima coluna de $\bvV$, também o autovetor de $\Corr{\bvs}$ referente ao autovalor $\lambda_n$. Assume-se que $\Corr{\bvs}$ seja uma matriz diagonalizável, tal que a auto-decomposição seja possível.

É possível reescrever a \cref{eq:sec5:ssp_combinacao-linear} utilizando a auto-decomposição de $\Corr{\bvs}$, levando a
\begin{equation}
	\bvs[t] = \bvs[b] + \bvV \bva
\end{equation}
em que $\bva$ é um (novo) vetor de coeficientes. Por fim, toma-se que apenas $G$ dentre os $N_{\z}$ autovetores são suficientes para reconstruir os SSP com acurácia. Ou seja,
\begin{equation}
	\bvs[t] = \bvs[b] + \bvV{G} \bva{G}
	\label{eq:sec5:modelagem_eof_truncado}
\end{equation}
em que $\bvV{G} \in \re^{\sz{N_{\z}}{1}}$ é composto pelas $G$ primeiras colunas de $\bvV$ --- essas, referentes aos $G$ maiores autovalores (em módulo).

\subsubsection*{Estimação do SSP}
\textsl{\footnotesize{} Nota: no artigo que inspirou essa formulação, foi utilizado um algoritmo genético para estimação das velocidades; aqui será usada a formulação por minimização do gradiente para estimação das vagarosidades.}

Seguiremos uma estrutura semelhante à da TTT, com minimização do erro de tempo de trânsito entre o modelo e o medido. Ou seja, o problema-base é
\begin{equation}
	E(\bva{i+1}) = \norm{\bvR{i} \bvs(\bva{i+1}) - \bvt{\obs}}^2
\end{equation}
com $\bvR{i}$ sendo a $i$-ésima solução da equação Eikonal; aqui toma-se que $\bva{G;i} \equiv \bva{i}$, por simplicidade da notação. Sabendo que $\bvs(\bva{i}) = \bvs[b] + \bvV{G} \bva{i}$, então a solução para esse problema será
\begin{equation}
	\bva{i+1} = \inv*{\bvV[T]{G} \bvR[T]{i} \bvR{i} \bvV{G}} \bvV[T]{G} \bvR[T]{i} \pts{ \bvt{\obs} - \bvR{i} \bvs[b] }
\end{equation}

\subsubsection*{Estimação do SSF}
A formulação acima presume que a estimação é uni-dimensional. Portanto, é necessário expandi-la para o caso multi-dimensional, com algumas adaptações.

Assumamos $N_{\x}$ colunas, cada uma com um SSP $\bvs{i,n} \in \re^{\sz{N_{\z}}{1}}$ ($0 \leq n < N_{\x}$). Disso, o SSF em forma matricial $\bvS{i} \in \re^{\sz{N_{\z}}{N_{\x}}}$ é dado por
\begin{equation}
	\bvS{i} = \bts{\bvs{i,0}~|~\bvs{i,1}~|~\cdots~|~\bvs{i,N_{\x}-1} }
\end{equation}
com $\bvs{i} = \rvec(\bvS{i})$ sendo a vetorização por linhas de $\bvS{i}$. Como cada SSP $\bvs{i,n}$ pode ser dado pela \cref{eq:sec5:modelagem_eof_truncado}, então
\begin{equation}
	\bvS{i} = \bvS[b] + \bvV{G} \bvA{i}
\end{equation}
com $\bvS[b]$ sendo a concatenação horizontal de $N_{\x}$ cópias de $\bvs[b]$, e $\bvA{i}$ a concatenação horizontal de todos os $\bva{i,n}$ ($0 \leq n < N_{\x}$), representando os coeficientes a serem estimados. Pelas propriedades da vetorização por linhas, temos que
\begin{equation}
	\begin{split}
		\bvs{i}
		& = \rvec(\bvS{i}) \\
		& = \rvec(\bvS[b]) + \rvec(\bvV{G} \bvA{i}) \\
		& = \rvec(\bvS[b]) + \pts{\bvV{G} \otimes \Id{N_{\x}}} \rvec(\bvA{i})
	\end{split}
	\label{eq:sec5:estimacao_ssf_otimizada_basica}
\end{equation}
com $\otimes$ simbolizando o produto de Kronecker.

Denotando $\bvV[t]{G} = \pts{\bvV{G} \otimes \Id{N_{\x}}}$, $\bva[t]{i} = \rvec(\bvA{i})$, e $\bvs[t]{v} = \rvec(\bvS[b])$, o problema torna-se
\begin{equation}
	\bvs{i} = \bvs[t]{v} + \bvV[t]{G} \bva[t]{i}
\end{equation}

Aliando esse resultado à estrutura de minimização quadrática,
\begin{equation}
	E(\bva[t]{i+1}) = \norm{\bvR{i} \bvV[t]{G} \bva[t]{i+1} + \bvR{i} \bvs[t]{v} - \bvt{\obs} }^2
	\label{eq:sec5:funcao-custo_eof_para_ssf}
\end{equation}
cuja solução é
\begin{equation}
	\bva[t]{i+1} = \inv*{\bvV[tT]{G} \bvR[T]{i} \bvR{i} \bvV[t]{G}} \bvV[tT]{G} \bvR[T]{i} \pts{\bvt{\obs} - \bvR{i} \bvs[b]{v}}
\end{equation}

Notavelmente, o perfil $\bvs{i+1}$ pode ser obtido pela \cref{eq:sec5:estimacao_ssf_otimizada_basica}.

\subsubsection*{Regularização}

Duas formas de regularização podem ser consideradas aqui: sobre $\bva{i+1}$ (o vetor de coeficientes da EOF), ou sobre $\bvs{i+1}$ (o vetor de vagarosidades reconstruído). Por essas duas variáveis serem linearmente independentes, regularizar uma afeta diretamente a outra; contudo, a diferente notação permite uma melhor visualização da penalização adicionada.

\paragraph{a)} Na primeira regularização, usa-se que a presença do perfil-médio $\bvs[b]$ no modelo implica que os coeficientes de $\bva{i+1}$ serão pequenos. Portanto, a adição de um termo de regularização $\alpha^2 \norm{\bva[t]{i+1}}^2$ à \cref{eq:sec5:funcao-custo_eof_para_ssf} modifica a solução para
\begin{equation}
	\bva[t]{i+1} = \inv*{\bvV[tT]{G} \bvR[T]{i} \bvR{i} \bvV[t]{G} + \alpha^2 \bvI{GN_{\x}}} \bvV[tT]{G} \bvR[T]{i} \pts{\bvt{\obs} - \bvR{i} \bvs[b]{v}}
	\label{eq:sec5:metodo_eof:reg_peq-coef}
\end{equation}

\paragraph{b)} Na segunda regularização, deseja-se melhorar a suavidade do perfil. Disso, o termo de regularização adicionado será $\alpha^2 \norm{\bvD \bvs{i+1}}^2$, que leva à solução
\begin{equation}
	\bva[t]{i+1} = \inv*{ \bvV[tT]{G} \pts{\bvR[T]{i} \bvR{i} + \alpha^2 \bvD[T] \bvD} \bvV[t]{G} } \bvV[tT]{G} \pts{ \bvR[T]{i} \bts{\bvt{\obs} - \bvR{i} \bvs[b]{v}} - \alpha^2 \bvD[T] \bvD \bvs[b]{v} }
	\label{eq:sec5:metodo_eof:reg_suavidade}
\end{equation}
Em ambos os casos, dado que a regularização não afeta o SSF objetivo, o perfil $\bvs{i+1}$ é reconstruído também utilizando a \cref{eq:sec5:estimacao_ssf_otimizada_basica}.

Ambas regularizações possuem formato semelhante, mas buscam impactar a solução de maneira diferente: a regularização (a) penaliza a variância dos coeficientes dos EOFs; enquanto a regularização (b) penaliza a variância (entre células) do SSF reconstruído.